\documentclass{article}
\usepackage{amsmath}
\usepackage{amsthm}
\usepackage{graphicx}
\usepackage[hidelinks]{hyperref}
\graphicspath{{Images/}}
\usepackage[
    left=2.15cm,
    right=2.15cm,
    top=1.90cm,
    bottom=1.90cm
]{geometry}

\title{3,741: Sustainable Minds - Personality Insights for Effective Climate Action}
\author{Leroy Zahno}
\date{Autumn Semester 2026}

\makeatletter
\renewcommand{\maketitle}{%
  \begin{center}
    {\LARGE \@title \par}
    \vspace{1em}
    {\large \@author \par}
    \vspace{1em}
    {\large \@date \par}
  \end{center}
}
\makeatother

\newtheorem{theorem}{Theorem}

\theoremstyle{definition}
\newtheorem{definition}{Definition}

\theoremstyle{Solution}
\newtheorem{solution}{Solution}

\begin{document}

\maketitle

\subsection{Disclaimer}

This document was written on the basis of the lecture slides from AS2026, and the additional readings from the course fact sheet. I use definitions and theorems to format this document, however these terms may at times not be an accurate description of the subject matter. 

\tableofcontents

\section{Lecture 1: What is personality and how to measure it}

\begin{definition}[Psychology]
Psychology is a field of science that aims to describe, explain, predict, and change Behavior, Cognition, and Emotion.
\end{definition}

\begin{definition}[Sub-fields of Psychology]
Clinical psychologists diagnose and threat mental health disorders. Biological Psychologists study the influence of the brain on behaviour ad mental functions. Social psychology studies the influence of social interactions on individual's thoughts, feelings and behaviours. Developmental psychologists study psychological changes across the lifespan. Educational psychology studies how people learn and proposes best practises for teaching. Organisational Psychology studies workplace behaviour, employee performance and well-being. Personality psychology studies individual differences in characteristic patterns of thinking, feeling, and behaving.
\end{definition}

\begin{definition}[Personality trait]
A personality trait refers to differences among individuals in a typical tendency to behave, think, or feel in some conceptually related ways, across a variety of relevant situations and across some fairly long period of time. 
\end{definition}

\begin{definition}[LOTS of data]
This is an abbreviation for Life records, Observational data, Test data, and Self-reported data. 
\end{definition}

\begin{definition}[Quality of measurement]
We say that data is sufficiently qualitative if it satisfies the conditions of objectivity, reliability and validity.
\end{definition}

\begin{definition}[Types of Reliability Measurement]
We commonly use two types of measurement to evaluate reliability of a research project. \textit{Internal consistency} checks if different items in the test correlate with one another as one wold expect if each item is a reflection of a common psychological construct. \textit{Test-retest reliability} checks if test scores remain the same or highly similar if the test is taken at two different points in time.
\end{definition}

\begin{definition}[Types of Validity Measurement]
For a test to be useful in the development and testing of personality theory it must have \textit{construct validity, or content validity}. Thus it must be shown that the test relates systematically to some external criterion. For example if a test measures the tendency to experience anxiety, we would use theoretical ideas about anxiety to choose external criteria, such as an index of anxious arousal. Then we would show that the test correlates with the external criterion.

We might also involve comparisons of two groups of people. For instance, people with clinical diagnoses of an anxiety disorder should get higher scores on the anxiety test than people who have not been diagnosed.

One should also be able to demonstrate \textit{discriminant validity}. This means that it should be empirically distinct from other tests that already exist. If we see that our test already highly correlates with existing tests of "worrying tendencies", the test is of little value.
\end{definition}

\begin{definition}[Structured Personality Inventory]
Those are questionnaires that are used for obtaining self- and observer reports of personality. They are called structured because they give the respondents a predetermined set of options for responding to the items. Otherwise they would be unstructured. Psychologists make roughly half of the items of a scale reverse-coded ore negatively keyed. This means questions are asked in a way to that expressing disagreement leads to higher scores on a given trait. Agreeing to "I like going out with friends" is equivalent to disagreeing to "I like spending time alone". This helps balance out the \textit{social desirability-} and \textit{acquiescence bias}.
\end{definition}

\section{Lecture 2: Models of personality}

\begin{definition}[The four humors]
According to Hippocrates, personality depends on equilibrium of the four fluids of the body. Those are the yellow bile (choleric), the black bile (melancholic), blood (sanguine), and phlegm (phlegmatic). However this theory does not hold empirically.
\end{definition}

\begin{definition}[The psychoanalytic model]
It suggests that personality is based on early childhood experiences. It claims that humans have the two basic drives to create (Eros) and to self-destruct (Thanatos). Additionally it is said that human personality consists of a conscious and subconscious part. The \textit{Id} is subconscious and seeks pleasure. The \textit{Superego} is partly conscious and represents cultural rules and authorities. The \textit{Ego} is mostly conscious and tries to find compromises between the Id and Superego. However this theory does not hold empirically either.
\end{definition}

\begin{definition}[Classical and operational conditioning]
Classical Conditioning (Pavlov, 1927): Shows how a \textit{conditioned stimulus (CS)} and \textit{unconditioned stimulus (US)} can be paired repeatedly, leading to a learned response to the conditioned stimulus alone. Operational Conditioning (Skinner, 1938): Shows that reinforced behaviour is shown more often and punished behaviour is shown less often. This basically means that reward and punishment have an influence on behaviour.
\end{definition}

\begin{definition}[Lexical hypothesis]
The lexical hypothesis states that people will want to talk about the personality traits that they view as having important consequences in their lives. As a result, people will inevitably invent some words to describe people who have high or low levels of these important traits. This process is called sedimentation. Allport \& Odbert found 17,953 terms to describe personality and behaviour.
\end{definition}

\begin{definition}[HEXACO]
This 6-factor model describes the following dimensions: Honesty-Humility, Emotionality, Extraversion, Agreeableness, Conscientiousness, and Openness to Experience.
\end{definition}

\begin{theorem}[Important characteristics of the OCEAN model]
Each trait exists on a continuum with an individual either scoring high or low on each one. These personality traits are relatively constant throughout adulthood. They can help predict various behaviours such as job performance or other life outcomes. These traits can be found across cultures, making them broadly applicable. Everyone has a mix of these traits and there are no good or bad combinations.
\end{theorem}

\begin{definition}[Myers-Briggs Type Indicators]
The MBTI consists of 4 dichotomous attributes, leading to \(2^4=16\) personality types. However there are several problems that make it inadequate for scientific usage: The model is overly simplistic, it has low predictive power of life outcomes, it shows poor test-retest reliability, and it ignores trait continuums.
\end{definition}

\begin{definition}[Factor analysis]
Factor analysis allows researchers to reduce many variables into a few unrelated groups of related variables.
\end{definition}

\section{Lecture 3: Development and stability of the Big Five}

\begin{definition}[Nature vs. Nurture]
Genetics provide the foundation for personality, determining broad predispositions for traits such as openness, conscientiousness, or extraversion. However, the environment (e.g., parenting style, education, relationships) plays a significant role in how these genetic predispositions are expressed and developed over time.
\end{definition}

\begin{definition}[Twin studies]
Some researchers wanted to learn about the impact of genes on personality. For that purpose they tested identical twins in the same environment, identical twins in different environments, twins with different genes in the same environment and twins with different genes in a different environment. They found out that identical twins tend to be twice as similar as unidentical twins. They claim that 50\% of that variance is due to the genetic differences between the unidentical twins.
\end{definition}

\begin{definition}[Types of Change and Stability]
\begin{enumerate}
	\item Stable: The personality of the individuals in a population does not change.
	\item Mean level change: Means that the average level of a trait changes in the population. But the relative standing of individuals remains the same (e.g. with old age all people become more conscientious). This means that we talk about personality aggregated by age.
	\item Rank order change: This happens when individual's traits change at different rates. This means their relative position in the population changes, compared to others.
	\item Mean and rank order change: The average level of some traits changes for the population at the same rate, however other traits vary in their rate of change.
\end{enumerate}
\end{definition}

\begin{theorem}[Typical mean level changes]
Conscientiousness and Agreeableness tend to increase with age. Assertive aspects of Extraversion tend to increase until middle age. Sociable aspects of Extraversion tend to decrease with age. Neuroticism tends to decrease until middle age. Openness tends to increase until adulthood. Overall people tend to develop a greater maturity throughout their life.
\end{theorem}

\begin{theorem}[Rank order stability]
This means, that a person who at the age of 20 is higher in conscientiousness than the average 20 year old has also a higher likelihood to be more conscientious at the age of 40 than the average 40 year old. This is confirmed by a test-retest reliability of ~.6 - .7 points, which means that there is a high correlation between the personality of an individual across two points in time (about 6-24 years apart). So once you reach adulthood, your personality is not going to change much, unless something specific happens here.
\end{theorem}

\begin{definition}[Personality change through Life Events]
Research shows that after childbirth people often experience a decrease in conscientiousness since parenting takes capacities necessary for leading a conscientious life. New relationships tend to increase conscientiousness and extraversion, while decreasing neuroticism. And finally, a first job tends to increase extraversion while decreasing neuroticism.
\end{definition}

\begin{definition}[Personality change through intervention]
Small to moderate personality change through interventions is possible. Some interventions to increase extraversion could be: Brainstorming lists of questions to ask other people; going to a public place where people mingle and saying hello to someone new; introducing yourself to a stranger; introducing yourself to a stranger and asking them at least two questions about themselves.
\end{definition}

\begin{definition}[Social investment hypothesis]
This is the idea that changes in personality between late adolescence and adulthood might reflect the greater sense of purpose or direction in life that many people develop after having become fully integrated into the adult world.
\end{definition}

\section{Lecture 4: Relevance of Personality}

\begin{theorem}[Relevance of Big Five for life outcomes]
The Big Five predict important life outcomes. Their predictive power is comparable to and sometimes better than socioeconomic predictors and other psychological predictors such as IQ. Openness tends to increase tolerance, aesthetic engagement and intellect, while decreasing conservatism and time spent watching TV. Conscientiousness tends to increase longevity, academic performance, sense of duty while decreasing substance use. Extraversion tends to increase popularity, performance in sales and life satisfaction. Agreeableness tends to increase relationship satisfaction, humility and honesty, while decreasing willingness to deceive others. Neuroticism tends to increase substance use, stress and spending time alone while decreasing relationship satisfaction, life satisfaction and performance in sales.
\end{theorem}

\begin{theorem}[Substance use]
People with low conscientiousness tend to use more substances, because they lack self-discipline and impulse control. They are also less likely to resist temptation. Pople with high neuroticism are also more likely to use substance because they may be anxious, moody or irritable and may feel tempted to temporarily control those negative emotions.
\end{theorem}

\begin{theorem}[Academic performance]
People with high conscientiousness tend to have better academic performances than others. This goes across all levels of education. This is because of high self-discipline, organisational skills and diligence. Thus they put forward more effort, efficiency and show a higher degree of attention to detail.
\end{theorem}

\begin{definition}[Psychological targeting]
Matz et al. (2017) show that by using persuasive communication we can encourage large groups to believe and act on our viewpoint. These efforts are more effective if the message is tailored to the audience's psychological characteristics. We can use this to do good things and encourage positive behaviour.
\end{definition}

\begin{definition}[Personality fit]
Florez et al. (2020) show that team productivity is higher when the team has similar Big Five profiles.
\end{definition}

\begin{definition}[Personality for Entrepreneurs]
Zhao et al. (2010) show that entrepreneurial success depends on entrepreneurial intentions to start a firm and the performance of the firm for its survival. People who are more likely to start a business and be successful with it are characterised by high openness, conscientiousness and extraversion. They are moderately agreeable and less neurotic than average.
\end{definition}

\begin{theorem}[Personality and peer relationships]
Newcomb et al. (1993) looked at differences between \textit{"average"} children and those who are identified either as \textit{popular, rejected, neglected, or controversial}. They found:
\begin{itemize}
	\item[(a):] Popular children tended to be low in aggression and high in sociability
	\item[(b):] Rejected children tended to be high in aggression and low in sociability.
	\item[(c):] Neglected children tended to be low in aggression and sociability.
	\item[(d):] Controversial children tended to be high in aggression and sociability.
\end{itemize}
However these results do not show the direction of cause. They don't show if being liked or disliked is the consequence or result of being sociable/unsociable.
\end{theorem}


\end{document}